%\needspace{4\baselineskip}
\color{uniblue}{\subsection{Защита от внутренних повреждений батареи статических конденсаторов}}
\color{black}

\begin{enumerate}[label=\arabic{section}.\arabic{subsection}.\arabic{enumi}, labelsep=4pt, leftmargin=0pt, itemindent=57pt, itemsep=0pt, parsep=0pt] % если что можно крутить itemsep


\item
Защита от внутренних повреждений (небалансная защита) предназначена для защиты батареи 
статических конденсаторов при замыкании ряда или пробое единичных конденсаторов. Защита 
контролирует токи небаланса, которые протекают в цепях проводников, соединяющих средние точки 
параллельных ветвей каждой фазы БСК. В случае равномерного повреждения конденсаторов 
параллельных ветвей в любой из фаз БСК возможен отказ данной защиты. 

\item
Естественный ток небаланса, протекающий в цепи проводника, соединяющего средние точки 
параллельных ветвей фазы БСК, может быть вызван: 

\begin{itemize}
    \item разбросом параметров конденсаторов в пределах допустимого значения, регламентируемого 
заводом-изготовителем; 
    \item изменением емкости под воздействием температуры окружающей среды (ввиду изменения 
свойств диэлектрика); 
    \item естественным старением конденсаторов. 
\end{itemize}

\item
Алгоритмом предусматривается предварительная компенсация естественного тока небаланса в 
фазах, выполненная в соответствии с формулами: 

\begin{equation}
\label{eq:INB}
I_{\textsl{НБ.х}} = \mid I_{\textsl{СТ.х}} - K_{\textsl{комп.х}} \cdot I_{\textsl{ВВ.х}} \mid
\end{equation}

\begin{eqexpl}[25mm]
\item{$х$} наименование контролируемой фазы (A, B, C); 
\item{$I_{\textsl{НБ.х}}$} значение тока небаланса с учетом компенсации; 
\item{$I_{\textsl{СТ.х}}$} значение тока небаланса, измеренное в цепи проводника, соединяющего средние точки 
параллельных ветвей фазы БСК; 
\item{$I_{\textsl{ВВ.х}}$} значение фазного тока, измеренное со стороны высоковольтных вводов БСК.  
\item{$K_{\textsl{комп.х}}$} коэффициент компенсации тока небаланса (индивидуальный для каждой фазы). 
\end{eqexpl}

\item
Защита содержит сигнальную и отключающие ступени. Каждая из ступеней имеет регулируемое 
значение по току срабатывания, задаваемое параметром \textbf{«Iсигн ЗВП»} и \textbf{«Iоткл ЗВП»} для сигнальной и 
отключающей ступени соответственно. Сигнальная ступень с задержкой по времени \textbf{«Тсигн ЗВП»} 
действует на сигнализацию, отключающая ступень с задержкой по времени \textbf{«Тоткл ЗВП»} действует на 
отключение выключателя БСК. 

\begin{figure}[H]
\centering
\includegraphics[width=1\textwidth,height=1\textheight,keepaspectratio]{\fbpath/Защита от повреждений батареи/img/ZVP.pdf}
\captionof{figure}{Функциональный блок контроля цепей тока}\label{pict:ZVP}
\end{figure}


\small
\setlength{\extrarowheight}{0.05cm} % высота ячеек в таблице
\begin{longtable}
{|>{\centering\arraybackslash}m{4.5cm}|>{\centering\arraybackslash}m{4.5cm}|>{\centering\arraybackslash}m{3cm}|>{\centering\arraybackslash}m{3.5cm}|}
\caption{Перечень уставочных параметров ЗВП\hfill\vspace{-0.5\baselineskip}}\label{ksn:tbl2} \\ 
\hhline{|-|-|-|-|}
\rowcolor{gray!30}
Наименование уставки & Диапазон & Шаг &  Значение по умолчанию \\ 
\hhline{|-|-|-|-|}
\endfirsthead

\caption*{\hspace{3pt}\emph{Продолжение таблицы \ref{ksn:tbl2}\hfill\vspace{-0.5\baselineskip}}} \\ % сделано по ГОСТ 2.105 п.6.8.7
\hhline{|-|-|-|-|}
\rowcolor{gray!30}
Наименование уставки & Диапазон & Шаг &  Значение по умолчанию \\  
\hhline{|-|-|-|-|}
\endhead
\hhline{|-|-|-|-|}

\endfoot
\endlastfoot

%%%%%%%%%%%%содержимое%%%%%%%%%%%%%%

\end{longtable}
\setlength{\extrarowheight}{0.0cm} % возврат параметра высота ячеек в таблице
\normalsize





















\end{enumerate} 